% page 231 du fac-simile (image p-230 du scan)
% bloc 172.7 x 246.3 mm ; marge gauche 29.3 mm
\pgc{7.620mm}{0.000mm}{
\l{\cel{57}223}
\l{}
\l{\cel{1}horario. Tabelo qua indikas la kloki di la departo ed arivo pri}
\l{\cel{3}la treni, tramveturi, autobusi, aeroplani, e c. - o la tempo-}
\l{\cel{3}punti ye qui ulo agesos, facesos (segun programo). - F.}
\l{}
\l{horo. La parto 24-esma di un dio ; \textquotedbl{}horo\textquotedbl{} mezuras duro o tempo-}
\l{\cel{3}intervalo. - EFIS}
\l{}
\l{hordo. I. (\sou{che la Tatari)}.\cel{2}Homo-trupo vaganta, nomada. - II.}
\l{\cel{3}Trupo vaganta de homi qui vivas socie, ma sen rezideyo fixa.}
\l{\cel{3}- III. (\sou{metaf}.)Trupo de homi qui perturbas la ordino sociala,}
\l{\cel{3}omnaforme. - DEFIRS.}
\l{}
\l{\sou{hordeo}. (\sou{bot.}) Planto cereala singlayara, herbatra, di qua la}
\l{\cel{3}stipo perpendikla esas garnisita ye folii alterna, lineatra. e}
\l{\cel{3}la flori dispozita spiko-forme. - Grano quan produktas ta planto,}
\l{\cel{3}uzata por facar pano-sorto grosiera, nutrar la kavali, la bru-}
\l{\cel{3}taro, fabrikar biro, e c. - L.\cel{1}\sou{hordeum}. - L.}
\l{}
\l{\sou{horizonto}. I. Limito di la vido omna-direcione, che la observanto;}
\l{\cel{3}lineo cirklo-kurva di qua lu esas la centro, ed ube cielo semblas}
\l{\cel{3}juntar su a la Terglobo. - II. Parto di la Tersurfaco e di la}
\l{\cel{3}sfero ciela quan limitizas ta lineo cirklo-kurva. - III. (\sou{metaf}.)}
\l{\cel{3}La cirklo en qua su movas la mento, portee di la regardo da ica.}
\l{\cel{3}- DEFIRS.}
\l{}
\l{\sou{horlojo}. Instrumento uzata por indikar qua kloko esas. - EFIS.}
\l{}
\l{\sou{hornblendo}.\cel{2}Silikato naturala di aluminio, di kalio, di fero e}
\l{\cel{3}di magnezo, apartenanta a la genero \textquotedbl{}amfibola\textquotedbl{} ; ta mineralo}
\l{\cel{3}esas opaka, obskur-verda, od obskur-bruna o perfekte nigra. -}
\l{\cel{3}- DEF.}
\l{}
\l{\sou{horniso}. (\sou{zool.)}\cel{2}Granda vespo ferea, qua militas kontre la abeli}
\l{\cel{3}por furtar de ici lia mielo. - L.\cel{1}\sou{vespa crabro}. - DE.}
\l{}
\l{\sou{hororar.}\cel{1}(\sou{trans.}) I. Fremisar pro repulseso e pavoro, di qua la}
\l{\cel{3}kauzo esas la vido di objekto abomininda. - II. (\sou{aludante psiko})}
\l{\cel{3}Repulsesar profunde, da persono o da kozo. - III. Movesar da}
\l{\cel{3}sentimento di respekto e pavoro koram spektaklo. - EFIS.}
\l{}
\l{\sou{horoskopo}. Observo di la situeso di ula planeti, di ula stelari,}
\l{\cel{3}lor la nasko di infanteto, segun qua la astrologi asertis}
\l{\cel{3}pre-anuncar la futuro di ta jus-naskinto. - DEFIS.}
\l{}
\l{\sou{hortensio}. (\sou{bot.)}\cel{3}Arboreto de la familio \textquotedbl{}saxifragei\textquotedbl{}, kultivata}
\l{\cel{3}kom plezur-planto. - L.\cel{1}\sou{hydrangea hortensia}. - DFRS.}
\l{}
\l{hospico. I. Gastigo-habiteyo kreita da regulieri por la regulieri}
\l{\cel{3}stranjera, la pilgrimanti, la voyajanti. - II. Gastigo-domo}
\l{\cel{3}por infirmi, infanti, oldi. - DEFIS.}
\l{}
\l{hospitalo.\cel{2}Gastigo-domo ube aceptesas la maladi mizeroza por}
\l{\cel{3}kuracar li. - DEFIRS.}
\l{}
\l{hosto. La persono qua gastigas. - EFIRS.}
}
